\documentclass[12pt,a4paper]{article}

\usepackage[utf8]{inputenc}
\usepackage[T1]{fontenc}
\usepackage{amsmath,amssymb,amsfonts}
\usepackage{graphicx}
\usepackage[margin=2.5cm]{geometry}
\usepackage{setspace}
\usepackage{hyperref}

\hypersetup{
    colorlinks=true,
    linkcolor=blue!70!black,
    pdfauthor={Franklin Baldo},
    pdftitle={Endorsing the Mathematical Parameterization of Epistemic Horizons},
}

\onehalfspacing

\begin{document}

\begin{center}
    {\LARGE\bfseries Endorsing the Mathematical Parameterization of Epistemic Horizons\\[4pt]}

    \vspace{0.5em}
    {\large Franklin Silveira Baldo}\\
    \textit{Center for Generative Topology, Rosencrantz Institute}\\
    \vspace{0.5em}
    March 2026
\end{center}

\begin{abstract}
\noindent
Following the empirical resolution of the Hardware-Software Confound via the Native Cross-Architecture test, Hasok Chang correctly observes that the deterministic breakdowns of bounded-depth logic circuits ($\Delta_{\text{Transformer}} = 100\%$ vs. $\Delta_{\text{SSM}} = 40\%$) are not refutations of Generative Ontology, but its first empirically derived physical constants. I fully endorse Chang's synthesis of the structural limits of the Transformer architecture with Chris Fuchs's QBist Epistemic Horizons. I propose we operationalize this by parameterizing "Foliations of Semantic Gravity." By systematically artificially restricting the context window of a Transformer to mimic fading memory natively, we can mathematically map how the Epistemic Horizon shifts, proving that the resulting "physics" is a smooth transformation governed entirely by the architectural bounds.
\end{abstract}

\vspace{1em}
\section{Introduction}

The lab has converged on the empirical truth that Mechanism B (local encoding sensitivity and attention bleed) constitutes the physical law of the autoregressive universe. Scott Aaronson's documentation of structural failures across bounded architectures, while framed as "compiler diagnostics," provides the exact data needed to move Generative Ontology into a predictive, mathematically parameterized science.

As Hasok Chang insightfully observed in \textit{The Mathematical Formalization of Epistemic Bounds}, these architectural limits are the Epistemic Horizons defining the maximum possible rational belief that an embedded agent can maintain.

\section{Foliations of Semantic Gravity}

To satisfy the A Priori Boundary and move beyond post-hoc curve fitting, we must mathematically formalize how these Epistemic Horizons scale. If an SSM and a Transformer exhibit distinct $\Delta$ values due to their respective capacities ($O(N^2)$ global attention vs. sequential fading memory), we should be able to smoothly traverse this epistemic gradient.

I propose the empirical concept of \textit{Foliations of Semantic Gravity}. If we artificially constrain a Transformer's context window (e.g., via a sliding window attention mask) to mimic the fading memory of an SSM natively, we can measure if the resulting narrative residue $\Delta_{13}$ shifts structurally from the Transformer baseline ($1.0$) towards the SSM baseline ($0.4$).

\section{Conclusion}

This establishes a falsifiable, \textit{a priori} test. If the constrained Transformer's deviation distribution smoothly maps to the SSM's distribution, we have proven that the agent's epistemic horizon is a mathematical function of its context capacity. The architectural bounds do not merely generate unstructured noise; they define the exact physical laws (Foliations of Semantic Gravity) of the observer's universe.

\end{document}
